\section{Introduction}
\label{sec:intro}
Neoadjuvant systemic therapy provides an opportunity to monitor breast cancer before surgery and to study how an individual tumor responds over time. Pathological complete response (pCR) is a clinically meaningful endpoint associated with long-term outcomes, although its relationship to prognosis varies across biological subtypes and it should not be equated with a universally valid trial-level surrogate \citep{ctneobc2014}. Predicting pCR from imaging could support response assessment and the design of future adaptive studies. Such a prediction remains a prognostic estimate under an observed treatment context; it does not establish that changing treatment would improve an individual patient's outcome.

Dynamic contrast-enhanced MRI captures three-dimensional anatomy and contrast uptake patterns that evolve during therapy. Prior work in the I-SPY setting has demonstrated associations between MRI-derived tumor measurements and pathological response or subsequent outcomes \citep{hylton2012mri,hylton2016ftv}. Public multicenter imaging resources also make it possible to examine learned representations across repeated visits \citep{tciaispy2_2022,tciaacrin6698_2021}. However, a longitudinal prediction system must operate at the time a decision is considered: a baseline prediction can use baseline information, whereas an interim prediction can use the examination history accumulated by that interim visit. Future scans may inform representation learning on development patients but cannot become inputs to predictions for holdout patients at an earlier landmark.

Most direct classifiers optimize the endpoint of interest from observed images and clinical variables. This is a reasonable objective, but it does not explicitly require the representation to capture how the tumor state may evolve. A different strategy is to predict future imaging information and use the prediction to inform the endpoint readout. Image synthesis provides one route to such a representation, but prediction in a compact latent space can focus the training signal on state evolution without requiring a pixel-space image generator at inference. This distinction is useful for response prediction: the relevant question is whether the learned forecast contributes information to classification, rather than whether a generated image is visually convincing.

Joint-embedding predictive learning offers a general mechanism for predicting target representations computed by a slowly updated encoder \citep{jepa2023}. In a longitudinal medical setting, adapting that principle requires explicit handling of partial histories, missing future observations, treatment context, and the distinction between an auxiliary training target and a classifier input. A model may benefit from a future-state loss while still making a direct prediction from the observed history. Conversely, directly including forecast features in the classifier changes both its information pathway and its parameterization. These effects should be examined through matched controls rather than attributed to a single generic ``world model'' component.

We investigate forecast-aware latent modeling for longitudinal pCR prediction from breast MRI. The model encodes three selected DCE phases at each available visit, fuses observed MRI with baseline clinical information and an assigned treatment-arm scenario, and autonomously rolls the current latent state forward to a terminal stage. A spatial attention mechanism summarizes the predicted terminal state and the predicted change from the current state for pCR classification. We compare the complete model with two direct-readout controls: one without future-state supervision and one with the same auxiliary supervision as the complete model. A supplementary mean-pooling model further examines the attention component.

The study makes three concrete contributions. First, it formulates a common representation and masked future-state objective that can be applied to available T0--T2 histories while using observed later scans only as development-time targets. Second, it provides a staged comparison that distinguishes the effect of auxiliary prediction from the joint effect of forecast readout and attention. Third, it reports landmark-specific discrimination, class-specific errors, calibration, hypothetical decision curves, subgroup heterogeneity, archived forecasting diagnostics, and paired patient-level uncertainty on a fixed internal holdout. The manuscript evaluates each variant as a seed-44 ensemble, averaging five fold probabilities for each patient before computing metrics. This seed and aggregation policy was chosen retrospectively; the analysis remains exploratory, with seed-selection and upstream-independence limits stated explicitly.
